\chapter{Introduction générale}
\label{chap:introduction}

\section{Contexte général}

La multiplication des cybermenaces et la sophistication croissante des
attaques informatiques ont conduit la majorité des organisations à se doter
de systèmes de détection et de supervision de sécurité, communément appelés
\textbf{SIEM} (\textit{Security Information and Event Management}). Ces
plateformes centralisent les journaux (\textit{logs}) issus des différents
équipements du système d'information — postes de travail, serveurs,
équipements réseau — et génèrent des \textbf{alertes} lorsqu'un comportement
suspect est détecté.

Dans ce projet, nous nous appuyons sur \textbf{Wazuh}, une plateforme SIEM/XDR
open source largement adoptée, qui combine un agent de collecte installé sur
les hôtes surveillés, un moteur de règles de détection, et un moteur de
recherche (\textit{Wazuh Indexer}, basé sur OpenSearch) qui stocke et indexe
les alertes générées.

Le problème central que rencontrent les équipes de sécurité opérationnelle
(\textbf{SOC}) n'est cependant pas la détection elle-même, mais le
\textbf{volume} d'alertes à traiter quotidiennement. Une étude régulièrement
citée dans l'industrie (rapports annuels de fournisseurs SIEM et enquêtes
SANS sur les opérations SOC) montre qu'un analyste peut être confronté à
plusieurs centaines, voire plusieurs milliers, d'alertes par jour, dont une
large proportion sont des \textbf{faux positifs}. Ce phénomène, connu sous le
nom de \textit{fatigue liée aux alertes} (\textit{alert fatigue}), dégrade la
capacité de détection réelle des menaces : les analystes finissent par
traiter les alertes de manière superficielle, voire par les ignorer.

\section{Problématique}

Comment réduire la charge de triage manuel des alertes de sécurité tout en
conservant — voire en améliorant — la qualité et la traçabilité des
décisions prises~? Plus précisément, ce projet cherche à répondre aux
questions suivantes~:

\begin{itemize}
  \item Peut-on automatiser la phase d'\textbf{investigation initiale} d'une
        alerte (collecte de contexte, recherche d'alertes corrélées,
        vérification de la réputation d'une adresse IP, etc.) à l'aide d'un
        agent d'intelligence artificielle~?
  \item Un tel agent peut-il produire un \textbf{verdict argumenté}
        (\textit{vrai positif}, \textit{faux positif}, ou \textit{à revoir})
        avec un niveau de confiance et une justification exploitable par un
        humain~?
  \item Comment garantir que cet agent reste \textbf{fiable, traçable et
        borné} (en coût, en temps, en accès aux données) alors qu'il
        s'appuie sur un modèle de langage dont le comportement n'est pas
        entièrement déterministe~?
  \item Comment restreindre l'accès de l'agent aux données sensibles
        (identités des utilisateurs, par exemple) sans complexifier son
        architecture~?
\end{itemize}

\section{Objectifs du projet}

L'objectif général du projet est de concevoir et développer
\textbf{\projectname}, une plateforme web qui~:

\begin{enumerate}
  \item collecte automatiquement les alertes générées par Wazuh
        (\textit{polling} périodique sur le Wazuh Indexer)~;
  \item détermine, à l'aide d'un \textbf{moteur de politiques}
        configurable, quelles alertes méritent l'ouverture d'une
        investigation~;
  \item confie chaque investigation à un \textbf{agent IA} (\textit{AI
        Investigator}) qui raisonne de manière itérative
        (\textit{boucle agentique}) : il consulte des outils (recherche
        d'alertes, historique des investigations, réputation d'IP,
        informations sur les actifs), accumule des preuves, puis rend un
        verdict argumenté~;
  \item enregistre \textbf{chaque action} de l'agent (outil appelé,
        raisonnement, résultat) pour permettre un audit humain complet~;
  \item propose un \textbf{assistant conversationnel} (chatbot) partageant
        les mêmes outils, pour interroger librement la plateforme en
        langage naturel~;
  \item sécurise l'accès à la plateforme par une
        \textbf{authentification} déléguée aux comptes Wazuh existants,
        évitant la duplication d'une base de mots de passe~;
  \item expose l'ensemble via une \textbf{interface web} (tableau de bord,
        liste et détail des investigations, gestion des actifs et des
        politiques, chat).
\end{enumerate}

\section{Méthodologie de gestion du projet}

Le projet a été mené selon une démarche \textbf{itérative et incrémentale},
proche des principes agiles~: chaque brique fonctionnelle (collecte
d'alertes, moteur de politiques, agent d'investigation, interface web,
assistant conversationnel, authentification) a été conçue, développée puis
validée indépendamment avant intégration, ce qui a permis d'obtenir à tout
moment une chaîne fonctionnelle de bout en bout, même partielle.

Le déroulement du projet a suivi les grandes phases imposées par le cadre
pédagogique du Projet de Fin d'Année~: définition du sujet, étude de
l'existant, conception, implémentation, présentation des résultats, et
rédaction du rapport. Le diagramme de Gantt ci-dessous synthétise cette
planification.

\begin{figure}[H]
\centering
\begin{ganttchart}[
    vgrid, hgrid,
    x unit=0.62cm,
    y unit chart=0.7cm,
    bar/.append style={fill=brandblue},
    bar height=0.55,
    group/.append style={fill=brandblue!60},
    milestone/.append style={fill=orange},
    title label font=\small\bfseries,
    bar label font=\small,
    group label font=\small\bfseries,
  ]{1}{16}
  \gantttitle{Semaines du projet (S1 -- S16)}{16} \\
  \gantttitlelist{1,...,16}{1} \\
  \ganttgroup{Phase 1 -- Définition du projet}{1}{2} \\
  \ganttbar{Cadrage sujet, contexte, objectifs}{1}{1} \\
  \ganttbar{Faisabilité \& méthodologie}{2}{2} \\
  \ganttmilestone{Jalon : validation du sujet}{2} \\
  \ganttgroup{Phase 2 -- Étude de l'existant}{3}{4} \\
  \ganttbar{État de l'art \& solutions existantes}{3}{4} \\
  \ganttgroup{Phase 3 -- Conception}{5}{7} \\
  \ganttbar{Architecture \& choix technologiques}{5}{5} \\
  \ganttbar{Modélisation UML \& base de données}{6}{7} \\
  \ganttgroup{Phase 4 -- Implémentation}{8}{13} \\
  \ganttbar{Collecte \& moteur de politiques}{8}{9} \\
  \ganttbar{Agent d'investigation IA}{9}{11} \\
  \ganttbar{Assistant conversationnel}{11}{12} \\
  \ganttbar{Authentification}{12}{12} \\
  \ganttbar{Interface web}{10}{13} \\
  \ganttgroup{Phase 5 -- Résultats \& validation}{14}{14} \\
  \ganttbar{Tests fonctionnels \& démonstration}{14}{14} \\
  \ganttgroup{Phase 6 -- Rédaction \& soutenance}{15}{16} \\
  \ganttbar{Rédaction du rapport}{15}{16} \\
  \ganttmilestone{Jalon : soutenance}{16}
\end{ganttchart}
\caption{Diagramme de Gantt du projet}
\label{fig:gantt}
\end{figure}

\section{Étude de faisabilité}

La faisabilité du projet a été évaluée selon trois axes~:

\begin{description}
  \item[Faisabilité technique] L'ensemble des briques technologiques
    retenues (Wazuh, Python/FastAPI, MySQL, Vue.js, API compatibles
    OpenAI pour les LLM) est \textbf{open source ou disponible en accès
    libre}, et documenté. Un environnement Wazuh de test est disponible ainsi
    qu'un accès à des API de LLM (Groq, compatible OpenAI) pour le
    prototypage rapide, avec la possibilité de migrer vers un modèle
    hébergé localement (Ollama) pour s'affranchir des contraintes de quota
    et de coût — cette portabilité est un objectif de conception du projet
    (voir \cref{chap:conception}).
  \item[Faisabilité temporelle] Le découpage en phases courtes
    (\cref{fig:gantt}) et la démarche incrémentale permettent de disposer
    en permanence d'une version fonctionnelle, réduisant le risque de ne
    pas aboutir à une démonstration exploitable en fin de projet.
  \item[Faisabilité humaine] Les compétences requises (développement
    backend Python, développement frontend JavaScript/Vue, notions de
    bases de données relationnelles, notions de sécurité informatique)
    correspondent au programme de la filière et sont réparties entre les
    membres du groupe.
\end{description}

\section{Plan du rapport}

Ce rapport est organisé en six chapitres. Le \textbf{chapitre 2} dresse un
état de l'art des SIEM, des approches d'automatisation du triage d'alertes
(SOAR) et des travaux récents sur l'usage d'agents fondés sur des modèles de
langage en cybersécurité. Le \textbf{chapitre 3} présente la conception de la
solution~: architecture globale, modélisation UML et modèle de données. Le
\textbf{chapitre 4} détaille l'implémentation~: environnement de travail,
traitement des données et description des modules réalisés. Le
\textbf{chapitre 5} présente les résultats obtenus et leur discussion.
Enfin, le \textbf{chapitre 6} conclut le rapport et ouvre sur les
perspectives d'évolution du projet.
